\subsubsection{Canonical Basis}
(also: Duquenne-Guigues basis or stem basis \cite{ganter_conceptual_2016}).
A formal context $(G, M, I)$ can also be defined via a set of implications $\mathcal{L}$ that is sound (i.e., each implication in $\mathcal{L}$ holds in $(G, M, I)$), complete (i.e., each implication that holds in $(G, M, I)$ follows from $\mathcal{L}$), and nonredundant (i.e., no implication in $\mathcal{L}$ follows from other implications in $\mathcal{L}$).
The canonical basis is $\{P \rightarrow P'' \smallsetminus P | P$ is pseudo-closed$\}$.
A finite subset $P \subseteq M$ is pseudo-closed if and only if (i) $P \neq P''$ (i.e., $P$ is not closed and in other words: $P$ is not a concept intent) and (ii) if $Q \subsetneq P$ is a pseudo-closed proper subset of $P$, then $Q'' \subset P$ (i.e., $P$ contains the closure of every pseudo-closed subset with fewer elements).
If (i) and (ii) holds true, $P$ is a pseudo-intent of $(G, M, I)$.
The canonical basis is nonredundant, i.e., no implication in the canonical basis is dispensable.
It is not the only implication basis, but there is no implication basis with fewer implications \cite{ganter_conceptual_2016}.